\documentclass[12pt]{article}

\usepackage{graphicx}
\usepackage{amssymb}
\usepackage{amsmath}
\usepackage[margin=1in]{geometry}
\usepackage{fancyhdr}
\usepackage{enumerate}
\usepackage[shortlabels]{enumitem}

\pagestyle{fancy}
\fancyhead[l]{Li Yifeng}
\fancyhead[c]{Homework \#1}
\fancyhead[r]{\today}
\fancyfoot[c]{\thepage}
\renewcommand{\headrulewidth}{0.2pt}
\setlength{\headheight}{15pt}

\begin{document}
	
	\section*{Exercise 2}
	
	\noindent Every two-element set of connectives containing $\lnot$ and one of $\{\land, \lor, \to\}$ is a functionally complete subset of (or in other words can be used to define) $\{\lnot, \land, \lor, \to, \leftrightarrow\}$.
	
	\bigskip
	
	\begin{enumerate}[start=1,label={\bfseries Part \arabic*:},leftmargin=0in]
		\bigskip\item Define $\land$, $\to$, and $\leftrightarrow$ using only $\lor$ and $\lnot$.
		
		\subsection*{Answer}
		
			\[\boxed{
				\begin{aligned}
					p \land	q			& := \lnot ( \lnot p \lor \lnot q), \\
					p \to q				& := \lnot p \lor q, \\
					p \leftrightarrow q	& := \lnot ( \lnot p \lor \lnot q) \lor\lnot (p \lor q)
				\end{aligned}
			}\]
			
		\bigskip\item Define exclusive disjunction ($\text{XOR}$, $\veebar$) using only $\land$ and $\lnot$.
		
		\subsection*{Answer}
		
			\[\boxed{
				p \veebar q := \lnot ( \lnot p \land \lnot q) \land \lnot (p \land q)
			}\]
			
		\bigskip\item Actually, there exists two binary operators (dual to each other), $\text{NAND}$ and $\text{NOR}$, that are complete by itself. The $\text{NAND}$ operators means ”not both” and can be defined as follows: $p \uparrow q \leftrightarrow \lnot(p \land q)$. The $\text{NOR}$ operators means ”neither nor” and can be defined as follows: $p \downarrow q \leftrightarrow \lnot(p \lor q)$. Pick one of the operators and show that it is enough to define all the usual operators $\{\lnot, \land, \lor, \to, \leftrightarrow\}$.
		
		\subsection*{Answer}
		
			\[\boxed{
				\begin{aligned}
					\lnot p				& := p \uparrow p, \\
					p \land q			& := (p \uparrow q) \uparrow (p \uparrow q), \\
					p \lor q			& := (p \uparrow p) \uparrow (q \uparrow q), \\
					p \to q				& := p \uparrow (q \uparrow q), \\
					p \leftrightarrow q	& := \Big( p \uparrow (q \uparrow q) \Big) \uparrow \Big( p \uparrow (q \uparrow q) \Big) \uparrow \Big( q \uparrow (p \uparrow p) \Big) \uparrow \Big( q \uparrow (p \uparrow p) \Big)
				\end{aligned}
			}\]
			
	\end{enumerate}
	
\end{document}
